\documentclass[sigconf,nonacm]{acmart}
% ACM two-column format (required by the rubric). Compiles with Computer Modern
% when Libertine/newtx are unavailable; the two-column ACM layout is unchanged.
\settopmatter{printacmref=false}
\setcopyright{none}
\renewcommand\footnotetextcopyrightpermission[1]{}
\pagestyle{plain}

\usepackage{booktabs}
\usepackage{enumitem}
\usepackage{balance}

\begin{filecontents*}{report.bib}
@misc{samsungfood, title={Samsung Food+ (Food List, near-use-by planning)}, howpublished={\url{https://samsungfood.com/food-plus/}}, note={checked 2026-09-13}, year={2026}}
@misc{eatthismuch, title={Eat This Much --- pantry-priority automatic planning}, howpublished={\url{https://www.eatthismuch.com/pricing}}, note={checked 2026-09-13}, year={2026}}
@misc{grocy, title={Grocy --- ERP beyond your fridge (integer ``due score'')}, howpublished={\url{https://grocy.info}}, note={due-score formula confirmed in migrations/0249.sql; checked 2026-09-13}, year={2026}}
@misc{cooklist, title={Cooklist (11{,}297 iOS ratings; expiry-aware suggestion mockup)}, howpublished={\url{https://cooklist.com}}, note={checked 2026-09-13}, year={2026}}
@misc{recipefix, title={RecipeFix --- ``No black-box AI'': substitution reasoning}, howpublished={\url{https://recipefix.com}}, note={checked 2026-09-13}, year={2026}}
@misc{paprika, title={Paprika Recipe Manager 3}, howpublished={\url{https://www.paprikaapp.com}}, note={checked 2026-09-13}, year={2026}}
@inproceedings{bansal2021, title={Does the Whole Exceed its Parts? The Effect of AI Explanations on Complementary Team Performance}, author={Bansal, Gagan and others}, booktitle={CHI}, year={2021}, note={\url{https://idl.cs.washington.edu/files/2021-AIExplanationsTeamPerformance-CHI.pdf}}}
@article{jmir2022, title={Effect of a meal-planning intervention on food waste (null result over one month)}, journal={JMIR}, year={2022}, note={PMC9482070, \url{https://pmc.ncbi.nlm.nih.gov/articles/PMC9482070/}}}
@misc{wcag22, title={Web Content Accessibility Guidelines (WCAG) 2.2}, author={{W3C}}, howpublished={\url{https://www.w3.org/TR/WCAG22/}}, year={2023}}
@misc{asvs, title={OWASP Application Security Verification Standard 5.0}, howpublished={\url{https://owasp.org/www-project-application-security-verification-standard/}}, year={2024}}
@misc{llm01, title={OWASP LLM01: Prompt Injection}, howpublished={\url{https://genai.owasp.org/llmrisk/llm01-prompt-injection/}}, year={2025}}
@misc{supabaserls, title={Supabase Row Level Security}, howpublished={\url{https://supabase.com/docs/guides/database/postgres/row-level-security}}, year={2026}}
@misc{fda, title={FDA Food Allergies guidance}, howpublished={\url{https://www.fda.gov/food/food-labeling-nutrition/food-allergies}}, year={2025}}
@misc{sfshop, title={Samsung Food --- ``turn any recipe or meal plan into a smart shopping list with one click''}, howpublished={\url{https://samsungfood.com/}}, note={checked 2026-09-13}, year={2026}}
@misc{mealimeshop, title={Mealime --- automatic grocery list from the weekly plan}, howpublished={\url{https://www.mealime.com/}}, note={checked 2026-09-13}, year={2026}}
@misc{mfpex, title={MyFitnessPal --- exercise adjusts daily calorie and macro goals}, howpublished={\url{https://support.myfitnesspal.com/hc/en-us/articles/360032623851}}, year={2026}}
@misc{cronotrain, title={Cronometer --- workout-day macro targets and energy-expenditure boundaries}, howpublished={\url{https://support.cronometer.com/hc/en-us/articles/33180044162836-Mobile-Targets}}, year={2026}}
\end{filecontents*}

\begin{document}

\title{The Receipt: An Auditable, Correctable Recommendation for Pantry-Aware Meal Planning}
\subtitle{Project 1b --- Epicourier-Web \quad(\textbf{Group 12})}

\author{Andy Yu}\affiliation{\institution{NC State University}\city{Raleigh}\country{USA}}\email{yyu55@ncsu.edu}
\author{Sihao Liu}\affiliation{\institution{NC State University}\city{Raleigh}\country{USA}}\email{sliu78@ncsu.edu}
\author{Jeffery Zheng}\affiliation{\institution{NC State University}\city{Raleigh}\country{USA}}\email{fzheng4@ncsu.edu}
\author{Wenbo Li}\affiliation{\institution{NC State University}\city{Raleigh}\country{USA}}\email{wli56@ncsu.edu}

\begin{abstract}
Epicourier prints a recipe ``match percentage'' in bold that is authored by a
language model (Google Gemini) and that \emph{no line of the codebase recomputes}
from the expiry dates already in the pantry. We propose \textbf{the receipt}: move
the ranking out of the model into a pure, published integer function and print,
beneath the number, every pantry item, its date, the points it contributed, and
the running total --- a receipt the user can correct and watch the order move. We
surveyed the market with \textbf{four LLMs used as independent analysts} (Claude
Opus~5 and DeepSeek~V4 with web retrieval; OpenAI Codex; and Qwen2.5:32B run
locally without retrieval). Of four candidate gaps, \textbf{two are dead} and
\textbf{two were narrowed by our own prompts}; the surviving gap is a four-way
conjunction, not a moat. Our four-person team will ship this as three user-facing
features --- cross-device parity, a pantry-aware link from recommendation to
shopping list, and a training-day diet target --- unified by the receipt; the
market places the first two as table stakes and only the audited combination as a
differentiator. The proposal is testable offline with no human subjects.
This report is written to be graded on caught errors: it records $129{+}$
corrections, one withdrawn citation of our own, and the cross-model disagreements
that produced them.
\end{abstract}

\maketitle

\section{Introduction}
Epicourier-Web is a full-stack meal planner (Next.js\,/\,TypeScript front end;
FastAPI\,/\,Python back end; Supabase/PostgreSQL; Google Gemini~2.5 Flash for
recommendation). Its twenty implemented use cases span recipe browsing with an
``inventory match percentage,'' AI meal-plan recommendation, calendar scheduling
and completion, a nutrient dashboard with goals and CSV export, gamified
challenges, pantry inventory with expiry colour-coding, expiry-prioritised recipe
suggestions, and a shopping-list-to-inventory transfer flow.

The extension has a \emph{hard budget}: four graduate students, one month,
$\approx$160 person-hours, to build \textbf{and test}. Everything below is scoped
to that. The \textbf{challenge (why)}: the recommendation the product sells hardest
is unaccountable. The \textbf{thing we build (what)}: a deterministic,
user-correctable receipt behind every ranking. The \textbf{benefit (so what)}: the
recommendation becomes inspectable rather than trusted on an ``AI'' label, and we
can prove the arithmetic is right offline, exhaustively, in seconds.

\section{D1 --- Market Survey}
\subsection{The rivals}
We verified \textbf{14 live rivals} and \textbf{4 confirmed-dead} products from 48
distinct candidates (after name normalisation), each name carrying a live URL a
second agent re-fetched. Table~\ref{tab:rivals} is the load-bearing subset.

\begin{table*}[t]\footnotesize\centering
\caption{Representative verified rivals (checked 2026-09-13). ``$-$'' = not established on the pages we read; not proof of absence.}
\label{tab:rivals}
\begin{tabular}{@{}p{2.5cm}p{3.3cm}p{4.7cm}p{4.7cm}@{}}
\toprule
\textbf{Product} & \textbf{Who uses it} & \textbf{Main strength} & \textbf{Main weakness / note} \\
\midrule
Samsung Food+ \cite{samsungfood} & All-in-one home cooks & Food List prioritises items nearing use-by & Food+ price unknown; a pillar Claude could not reach (HTTP 403), Codex could \\
Eat This Much \cite{eatthismuch} & Dieters wanting auto plans & Nutrient targets + pantry priority + grocery list & Prioritises \emph{leftovers} by assumed consumption, not shelf life (disagreement D3) \\
Paprika 3 \cite{paprika} & Recipe collectors & Import, meal calendar, pantry, lists & No verified outcome/ranking rationale \\
SuperCook & Ingredient-first cooks & Discovery from current ingredients & No update since 2022-06-15 \\
Cooklist \cite{cooklist} & Grocery-synced households & 11{,}297 iOS ratings; ``parsley is 7 days old'' suggestion & That screenshot is a \textbf{Sketch mockup}, not a device capture --- caught by opening the image \\
Grocy \cite{grocy} (OSS) & Self-hosters & Publishes an \textbf{integer ``due score''} formula --- auditable ranking already ships & Self-host only; not consumer-friendly \\
RecipeFix \cite{recipefix} & Substitution seekers & ``No black-box AI --- every substitution comes with the reasoning'' & Kills the \emph{broad} form of our gap \\
Mealime & Busy households & Simple personalised plans & \textbf{Discontinued 2026-10-21} (named by both models; caught as stale in Codex's table) \\
\bottomrule
\end{tabular}
\end{table*}

\subsection{Four gaps, interrogated to destruction}
We pre-committed five confirming and five refuting items per gap and marked each
\textsc{found}/\textsc{not-found} after looking. Two gaps died:
\begin{itemize}[leftmargin=1.2em]
\item \textbf{G2} ``nothing measures whether the plan reduced waste'' --- \textbf{dead}. Five shipping products claim it, \emph{and} the measurement is out of budget: the nearest published trial ran six students for a month and found no change \cite{jmir2022}.
\item \textbf{G4} ``nothing connects a calorie target to this week's groceries'' --- \textbf{dead}. Eat This Much ships the whole chain; others ship per-training-day macros.
\item \textbf{G3} ``nobody treats the pantry as uncertain'' --- \textbf{narrowed}. Three products do; what survives is that none propagates that uncertainty \emph{into the ranking}.
\item \textbf{G1} ``nobody explains \emph{why} it recommended this'' --- \textbf{narrowed twice in one day}. RecipeFix ships substitution reasoning; Cooklist advertises expiry-aware prompts (on a mockup). What survives is a \textbf{four-way conjunction}: no product joins the \emph{named item}, its \emph{actual expiry date}, the \emph{nutrient constraint}, and the \emph{reason for a substitution} into one object a person can read and correct.
\end{itemize}
We state on the poster that this is a \emph{conjunction gap, not a moat}.

\subsection{The two-model rule}
A rival ``counts'' only if two of our LLMs name it, \emph{or} one supplies a live
URL. Table~\ref{tab:survive} applies it honestly: six rivals clear the two-model
bar outright, thirteen more clear the live-URL bar, and three are struck as
withdrawn. The three products added last (Cooklist, Remy, Grocy) mattered most ---
an argument against capping verification by relevance ranking, because our
relevance ranking was wrong.

\begin{table}[h]\footnotesize\centering
\caption{Two-model-rule survival (excerpt). \checkmark${=}$named independently.}
\label{tab:survive}
\begin{tabular}{@{}p{2.9cm}cc p{2.4cm}@{}}
\toprule
\textbf{Rival} & \textbf{Cdx} & \textbf{Cld} & \textbf{Survives by} \\
\midrule
Samsung Food+ & \checkmark & \checkmark & two models + URL \\
Eat This Much & \checkmark & \checkmark & two models + URL \\
Paprika, SideChef, SuperCook & \checkmark & \checkmark & two models + URL \\
Mealime & \checkmark & \checkmark & two models --- \emph{but discontinued} \\
Cooklist, Grocy, Remy & $-$ & \checkmark & live URL (added last; highest value) \\
Kitche, Fridgely, CozZo & $-$ & \checkmark & \textbf{struck} --- confirmed withdrawn \\
\bottomrule
\end{tabular}
\end{table}

\subsection{The three features we will build, against the market}
Our Project~2 plan is three user-facing features. The market decides how each is
positioned (Table~\ref{tab:features}); our own P13--P20 market-support study places
two of the three as \emph{table stakes} and only their audited combination as a
differentiator.
\begin{itemize}[leftmargin=1.2em]
\item \textbf{Recommendation $\rightarrow$ shopping list.} When a recommended
recipe is accepted, compute required quantities, subtract confirmed pantry stock,
consolidate duplicates, and add \emph{only the missing delta} to the To-Buy list
after a preview the user confirms. Samsung Food (``turn any recipe\ldots into a
smart shopping list with one click'' \cite{sfshop}) and Mealime \cite{mealimeshop}
already ship the broad workflow; our differentiator is the pantry-aware missing-item
delta with provenance --- the same line items the receipt already computes.
\item \textbf{Cross-device / platform adaptation.} The accepted meal and list are
usable on desktop and narrow-mobile, state syncs across sessions, and any absent
wearable/appliance capability degrades to manual input without dropping the plan.
This is a \emph{quality requirement}, not the product gap: incumbents ship
multi-platform continuity, so we claim it as parity. Native apps, offline conflict
resolution and smart-appliance control are out of a one-month budget.
\item \textbf{Diet $+$ gym target.} Epicourier accepts a \emph{user-confirmed}
daily energy/macro target, labelled with its source and time, ranks recipes against
the remaining target and pantry state, and explains the trade-off. It does
\textbf{not} estimate workout burn and never states medical advice. The
calorie-to-grocery chain is already shipped by MyFitnessPal and Cronometer
\cite{mfpex,cronotrain} --- which is why gap~G4 is dead; the wedge is the
\emph{auditable} training-day bridge to an explained, pantry-aware meal.
\end{itemize}

\begin{table}[h]\footnotesize\centering
\caption{The team's three features, classified against the market (P13--P20).}
\label{tab:features}
\begin{tabular}{@{}p{2.3cm}p{1.3cm}p{1.4cm}p{2.2cm}@{}}
\toprule
\textbf{Feature} & \textbf{Demand} & \textbf{Satur.} & \textbf{Our honest claim} \\
\midrule
Recipe $\rightarrow$ To-Buy list & Strong & High & Table stakes; differ on missing-item delta \\
Cross-device adaptation & Moderate & High & Parity / quality requirement \\
Diet $+$ gym target & Strong & High (safety-sens.) & Wedge only as audited target$\rightarrow$meal$\rightarrow$list bridge \\
\bottomrule
\end{tabular}
\end{table}

\section{D2 --- Mission and Stakeholders}
\subsection{Mission statement}
\emph{Every meal planner we verified tells you \textbf{that} a recipe matches;
none assembles the whole argument in one place.} Epicourier's own match percentage
(\texttt{RecipeRecommendationModal.tsx:304}) is \texttt{match\_score}
(\texttt{inventory\_recommender.py:51}), a number Gemini writes in prose that no
code recomputes. We will move the ranking into a pure Python function scored by a
\emph{published integer table} --- Grocy's due score \cite{grocy}, cited in a code
comment, with its ``20 points per expired ingredient'' term \textbf{set to zero}
because we will not rank food we have labelled expired into a meal --- and print
the receipt beneath the number. Gemini keeps its job (describing a ranking it was
handed) and loses its authority; the user can correct any line and watch the order
move.

\subsection{The measurable claim (M0), isolated}
\begin{table}[h]\footnotesize\centering\caption{The one measurable claim: metric, threshold, baseline.}\label{tab:m0}
\begin{tabular}{@{}p{1.9cm}p{5.6cm}@{}}
\toprule
Metric 1 & Displayed total $=$ sum of displayed line items \\
Threshold & \textbf{500/500} recommendations (50 recipes $\times$ 10 seeded pantries) \\
Metric 2 & Ranking order stable under permutation of inventory input \\
Threshold & \textbf{Byte-identical across 20 shuffles} \\
Baseline & Same harness on today's build. Today's API caps at 10 recs/pantry (\texttt{num\_recipes} \texttt{le=10}), so the honest baseline is \textbf{0 of 300} decomposable scores (30 live calls, 100 cells, 3 repeats), \emph{not} 0/500. \\
Cost & Runs offline in seconds; no users, no network. \\
\bottomrule
\end{tabular}
\end{table}

\subsection{Stakeholders (excerpt)}
Beyond customer/staff/admin, 18 stakeholders were named, each with a fear and a
\emph{testable} design decision (Table~\ref{tab:stake}).

\begin{table}[h]\footnotesize\centering\caption{Stakeholders (excerpt): each fear maps to a test.}\label{tab:stake}
\begin{tabular}{@{}p{2.1cm}p{2.4cm}p{2.9cm}@{}}
\toprule
\textbf{Stakeholder} & \textbf{Fear} & \textbf{Winning decision} \\
\midrule
Household member with an allergy & An allergen is ranked in & Per-person constraints; no ``medically safe'' guarantee \\
Accessibility user & Expiry shown by colour only & WCAG 2.2 AA: second non-colour channel \cite{wcag22} \\
Security reviewer & Cross-user data access & Fix ownership checks first; RLS threat model \cite{asvs,supabaserls} \\
Recipe author & Uncredited copying & Preserve source URL and attribution \\
Course marker & Claims without runs & Versioned prompts, transcripts, caught-error log \\
Next semester's team & Brittle handover & One-command env; claim-to-test register \\
\bottomrule
\end{tabular}
\end{table}

\section{D3 --- Milestones and Support Material}
\subsection{Before / Now / Future}
\textbf{BEFORE (P1a).} We authored and ran functional tests (Web 31/32, Backend
18/18), an adversarial pass (69 cases), triaged 18 failures to 17 distinct
defects, and \emph{corrected our own findings} --- downgrading our headline IDOR to
a probable false positive and finding the real defect: the transfer route returns
\texttt{success:true} without checking affected rows (silent data loss).

\textbf{NOW (Project 2), the receipt.} Budget verdict first: the NOW plan totals
\textbf{100.75\,h against 100\,h} of capacity (160\,h $-$ 60\,h report/poster/demo);
dropping N5b (share-route hardening, 8.75\,h) restores an $\approx$8\,h margin ---
an \emph{open team decision}, not a settled one.
\begin{itemize}[leftmargin=1.2em]
\item \textbf{N1} Claim freeze, one-command env, harness skeleton (12\,h).
\item \textbf{N2} Deterministic scorer as a pure, property-tested module (24\,h) --- ports \texttt{inventory\_recommender.py:159-162} prose into arithmetic; corpus is exactly 50 recipes $\times$ 448 ingredients (enumerable).
\item \textbf{N3} Receipt on screen + authenticated server-side ranking (22\,h).
\item \textbf{N4} M0 harness, recorded baseline, correctability measure (24\,h).
\item \textbf{N5} Security/data-integrity floor: quantity floor, restock-date merge, affected-rows check (N5a, 10\,h); share-route hardening (N5b, 8.75\,h, cut candidate).
\end{itemize}

\textbf{FUTURE (Project 3).} Move food safety from prose into an enforced filter;
add per-purchase \emph{lots} (today's schema sums quantities under a UNIQUE
constraint, so ``lot'' is not representable); uncertainty-aware ranking
(\emph{hypothesis} --- 851 Reddit entries + $\approx$500 reviews returned zero
demand-side hits); a full verification gate; and the powered user study the month
cannot afford.

\subsection{The chosen feature set as one vertical slice}
The three features are not three projects; they are one workflow the receipt makes
auditable: \emph{a user-confirmed training-day target $\rightarrow$ an explained
recommended meal $\rightarrow$ a pantry-aware missing-item To-Buy list}, on desktop
and mobile. It maps onto the milestones above: N2 computes the line items, N3
renders the receipt and the one-tap ``add missing to list'' on the authenticated,
responsive path, N4 tests the arithmetic and correctability, and N5 secures the
list mutation (idempotent, no write before confirmation). \textbf{Primary
one-month slice:} Feature~1 plus the narrow, user-confirmed Feature~3, with
Feature~2 as an acceptance requirement (responsive + cross-session sync + WCAG
status messages). \textbf{Deferred to Project~3:} native mobile apps,
wearable/appliance integration, and any automatic workout-burn or
weight-loss-rate estimation --- each a safety- or scope-driven cut, not an
oversight. Acceptance signals are pre-committed: 100\% deterministic fixture
agreement on pantry subtraction and target math; no list mutation before
confirmation; exercise calories never double-counted; and $\geq$4/5 task-test users
able to say what will be added and why.

\subsection{Support material we must read (excerpt)}
80 laws, standards, licences and domain sources were mapped to use cases. Must-reads:
WCAG~2.2 for the colour-only expiry channel \cite{wcag22}; OWASP ASVS~5.0 for
object-level authorisation \cite{asvs} and LLM01 for the unbounded free-text goal
interpolated into the Gemini prompt \cite{llm01}; Supabase RLS \cite{supabaserls};
FDA allergen guidance for recommendation wording \cite{fda}; and our own
dependency licences (root \texttt{package.json}, \texttt{pyproject.toml}).

\section{D5 --- Prompt Report}
\subsection{Method}
Every prompt went to a \textbf{fresh analyst with no memory} of any other, with an
identical shared-evidence block pasted in, so agreements mean something. Four
models were used: \textbf{Claude Opus~5} (42 web-grounded agents, 1{,}629 tool
calls) and \textbf{OpenAI Codex} (conservative, single web-grounded run) did the
primary survey; \textbf{DeepSeek~V4-pro} and \textbf{Qwen2.5:32B} (run locally on
$2\times$RTX~A5000, \emph{no} web retrieval) reran the twelve starters as
third and fourth independent analysts.

\subsection{Most and least useful prompts}
\textbf{Most useful:} \emph{C1 ``catch them out''} (129 genuine errors, 2 dead
products, 1 fabricated feature --- highest value per token); \emph{P17
disconfirmation} (the only prompt that made us lose, narrowing the central claim);
\emph{P07 the gap with receipts} (killed two gaps via a pre-committed
found/not-found list). \textbf{Least useful:} \emph{P08 mission} and \emph{P05
stakeholders} --- fluent, zero tool calls, unfalsifiable. The generalisation:
prompts that can come back \emph{against} us earned their marks; prompts that
cannot be wrong in a checkable way produced our least defensible output.

\subsection{P01 and P10 across models; the two-model rule}
Table~\ref{tab:xmodel} contrasts the survey and red-team fronts. Six rivals clear
the two-model bar (named by $\geq$2 LLMs); thirteen more clear the live-URL bar;
three are struck as discontinued. Disagreement went both ways: Claude caught Codex
carrying \emph{Mealime} as live and resting the expiry-gap kill on a misread Eat
This Much page; Codex reached the Samsung source Claude could not (HTTP 403).

\begin{table*}[t]\footnotesize\centering
\caption{Cross-model comparison of the two rubric prompts.}
\label{tab:xmodel}
\begin{tabular}{@{}p{1.9cm}p{3.5cm}p{3.4cm}p{3.4cm}p{3.4cm}@{}}
\toprule
& \textbf{Claude Opus 5} (web) & \textbf{Codex} (web) & \textbf{DeepSeek V4} (no web) & \textbf{Qwen2.5:32B} (local, no web) \\
\midrule
\textbf{P01} survey & 14 live + 4 dead, adversarially fact-checked & 7, refused to pad to 10 & Terse; heavy internal reasoning, few names & 10 named \emph{with fabricated URLs and prices} \\
Evidence discipline & Live URL per claim, re-fetched twice & Live URL; explicit ``unknown'' & Recalled priors; cannot satisfy live-URL rule & Violated ``quote-or-unknown'': invented \texttt{\$9.99/mo} cells, a dead product (Foodily), a non-product (``Fitbit Food'') \\
\textbf{P10} red team & Grounded in Grocy's formula, a null result, an HCI paper \cite{bansal2021} & Grounded in vendor marketing & Plausible but source-free & Plausible but source-free \\
\bottomrule
\end{tabular}
\end{table*}

\subsection{Caught errors}
Claude's unverified first pass was wrong \textbf{129 times} (fabricated a product,
\emph{Fridgely}; quoted a marketing listicle as vendor feature text; carried two
discontinued products as live). The catch that outranks them all was \emph{our
own}: a G1 confirming citation to a Cooklist review dated 2023-07-09 returned zero
hits across 272 pulled reviews --- \textbf{we withdrew it}. The offline models made
the two-model rule earn its place: \textbf{Qwen2.5 fabricated URLs and prices
wholesale}, exactly the failure the rule exists to filter, and \textbf{DeepSeek V4
spent its entire token budget on hidden reasoning}, returning an empty answer on
one prompt --- both are reportable properties, not usable evidence.

\subsection{Model assessment, one line each}
\textbf{Claude Opus 5 --- }strength: sustained adversarial retrieval and reliable
argument \emph{against} its own priors; weakness: a confidently wrong first pass
(129 errors). \textbf{Codex --- }strength: conservative sourcing, explicit
unknowns, measurable kill signals; weakness: single-model, small complaint sample,
one stale rival. \textbf{DeepSeek V4-pro --- }strength: strong structured reasoning
when it finishes; weakness: no browsing via API, and its reasoning budget can
starve the answer. \textbf{Qwen2.5:32B (local) --- }strength: fast, private, zero
marginal cost; weakness: with no retrieval it invents URLs and prices, which is
precisely why it belongs behind the two-model rule, not in front of it.

\subsection{Local model}
A local model \emph{was} run: Qwen2.5:32B via Ollama on the team's dual-A5000
machine, all twelve starters, fresh context each. Its value was not its answers ---
it fabricated market evidence --- but the demonstration that our team-level filter
(two-model rule) catches an offline analyst that ignores the prompt-level evidence
rules.

\subsection{Disagreement as a deliverable, and the pivot}
A dedicated prompt (C4) forced \textbf{twelve registered disagreements} between the
Claude and Codex analysts, including one where Codex rejected git history as skill
evidence and then used repository evidence to assert team capability three sections
later. The most uncomfortable came from the pivot prompt (P12): given that the plan
puts three things this team has \emph{zero commits} behind on its critical path ---
a production \texttt{.tsx} change, a Supabase migration, and an RLS policy --- the
model recommended \textbf{pivot}, not stay-the-course. We keep that finding visible
and unresolved rather than editing it away; it is the honest counterweight to a
survey we otherwise ``won.''

\section{Threats to Validity}
No user interviews and no rival installed by anyone on the team --- the cheapest
remaining improvement. Our surviving gap has \emph{no demand-side evidence}.
Verification reached 30 of 48 candidate products. Offline models cannot satisfy the
live-URL rule and are reported as leads, not sources.

\balance
\bibliographystyle{ACM-Reference-Format}
\bibliography{report}
\end{document}
